\scanpage{84}


Si on se donne (loc\textsuperscript{t.} ?) des bases \deleted{\unsure{de}} \inserted{$\{e_1\}$} de $L_1$, $\{e_2\}$ de $L_2$, donc $b_1,b_2$ s'interprèt\deleted{ent} comme des sections de $O$, alors la condition est

\[ (b_1,b_2,\underset{\substack{\Vert\\2\,\text{en}\\\text{l'occurrence}}}{\unsure{p|\gamma|\chi}})\cdot\underline{O}_X=\underline{O}_X \]

Cette condition est vérifiée si $p$ \inserted{(soit $2$ en l'occurrence)} inv. sur $X$. La condition b) aussi, comme on voit \deleted{en} la décomposition canonique

\[ E_i\simeq L_i\oplus O_i\qquad\text{avec }\sigma_i\mid O_i=+{\rm id}_{O_i} \]

d'où

\[ E_1\otimes E_2\simeq L_1\otimes L_2\oplus L_1\oplus L_2\oplus\underline O \]

et

\[ \sigma_1\otimes\sigma_2\quad\text{est}\quad\begin{cases}{\rm id}\text{ sur }L_1\otimes L_2\text{ et sur }O\\-{\rm id}\text{ sur }L_1\oplus L_2\end{cases} \]

et

\[ {\rm id}-\sigma_1\otimes\sigma_2\quad\text{est donc}\quad\begin{cases}0\text{ sur }L_1\otimes L_2\text{ et sur }O\\2\,{\rm id}\text{ sur }L_1\oplus L_2\end{cases} \]

donc

\diagram{\begin{tikzcd}[ampersand replacement=\&,column sep=small,row sep=large] (E_1\otimes E_2)_{/\sigma_1\otimes\sigma_2}\arrow[r,phantom,"\simeq"] \& (L_1\otimes L_2)\oplus O\arrow[r,"\sim"]\arrow[dr,"{x_1\otimes x_2\mapsto2x_1\otimes x_2}"',near end]\arrow[dr,bend left=18,"{\lambda\mapsto\lambda}"] \& \underbrace{E_1*E_2}_{\substack{\text{avec son}\\\text{scindage}\\\text{canonique}}}\\ \& \& L_1\otimes L_2\oplus O\end{tikzcd}}

Donc le seul pbl se pose en car. $2$, comme de juste. Je dis qu'\deleted{\unsure{alors}} alors (\unsure{Si on prend|Si on a} $L_1,L_2$ de prés. finie\unsure{\}|)}, p. ex. des Modules inversibles) la condition b) est \uline{conséquence} de a).
% slanted

\margin{\deleted{\unsure{On}} noter si on suppose que $L_1,L_2$ sont loc libres, alors dans $E=\unsure{\mathscr C_1|C_1}E_1\otimes E_2$ \unsure{$\mathscr C_i$} \unsure{$\operatorname{im}$} $(E_1\otimes E_2\longrightarrow E_1*E_2)$ la formation des $\unsure{\mathscr C_i}$ commute au changement de base $X_0\hookrightarrow X$.}

Il suffit le vérifier \deleted{\unsure{au ca}} si $X_{\!0}=\deleted{\unsure{X}}\unsure{X_0|X}$, de car. $2$ \textsuperscript{*} et il suffit de voir que si alors $T_1\in\Gamma L_1$
% archive sheet number: 83
